技术类书籍中常常包含作者的告诫，要求读者必须做练习与习题，这并不罕见. 我读到这类警告时总感到有些奇怪. 如果我不做练习与习题，会有什么坏事发生吗？当然不会. 我会节省一些时间，但代价是理解不够深入. 有时这值得. 有时则不值得.

那么这本书中什么值得做呢？我的建议是，你确实应该尝试做大部分练习，而你应该以\emph{不}做大部分习题为目标.

你应该做大部分练习，因为它们是检验你是否理解了材料的基本检查. 如果你不能相对轻松地解决一道练习，你很可能遗漏了某些根本性的东西. 当然，如果你偶尔在某道练习上卡住了，继续往下走就好——很可能只是你有一点小误解，或者也许是我某处表述得不好. 但如果大部分练习都做得很吃力，那么你大概需要重读一些前面的材料.

习题则是另一回事. 它们比练习更难，你很可能在某些习题上苦苦挣扎. 这很烦人，但当然，面对这种挫败感保持耐心，是真正理解和内化一个主题的唯一途径.

话虽如此，我并不建议把所有习题都做一遍. 更好的做法是找到你自己的项目. 也许你想用神经网络来分类你的音乐收藏. 或者预测股票价格. 或者别的什么. 但\emph{找到一个你在乎的项目}. 然后你就可以忽略书中的习题，或者仅仅把它们当作你自己项目工作的灵感来源. 在一个你在乎的项目上苦苦挣扎，会比做任意数量的规定习题教给你更多东西. 情感上的投入是达到精通的关键.

当然，你可能并没有想到这样的项目，至少一开始没有. 这没关系. 去做那些你有动力去做的习题. 并利用书中的材料来帮助你寻找有创意的个人项目的想法.
